% =============================================================================
% Section 11 results fragment: robustness and validation of the pillars, the Strategic Index,
% the allied picks and the PSI. Every number below is written to paper/results/robustness.json
% by scripts/paper/fig_robustness.R, which also draws the seven figures (robustness_*.pdf) and
% writes the six tables (paper/tables/robustness_*.tex). Inputs: the FINAL full rerun of
% 2026-10-12 (index_outputs.rds, theme tables, vintages, psi_*_tbl.rds), the October 2026 NIPO
% component (<scratch>/nipo_oct), the vertical-complementarity study (<scratch>/vertical) and the
% CM-FDI integration study (<scratch>/fixes7, inputs of 2026-10-11; not recomputed).
% Rank statistics are Spearman rank correlations. Cross-country rankings are taken within the
% technology-stage pairs of the eleven headline technologies over RANKABLE cells only
% (rankable == TRUE: sovereign, imputed weight share <= tau = 0.5).
% Labels defined here: tab:robustness_summary, tab:robustness_importance,
%   tab:robustness_rankability, tab:robustness_picks, tab:robustness_data,
%   tab:robustness_vertical (via \input of paper/tables/), eq:variance_share,
%   fig:robustness_variants, fig:robustness_vintage, fig:robustness_ranks, fig:robustness_loo,
%   fig:robustness_si, fig:robustness_picks_data, fig:robustness_psi_views.
% Table placement: tab:robustness_vertical is the appendix sensitivity table of the
%   vertical-complementarity study (user decision of 2026-10-12: paper only); it is input at the
%   end of this fragment and may be moved to an appendix by the integrating section.
% Section label: the integrating sections/11_robustness.tex must define \label{sec:robustness},
%   which Sections 1, 2, 4, 8, 9, 10 and Appendix D cite.
% External references: fig:strategic_sign, fig:strategic_screen (results/text/strategic.tex),
%   fig:method_rankability (results/text/method.tex), sec:network, sec:bilateral, eq:si,
%   eq:si_allied, eq:affine_constant, eq:scurve_logit, eq:scurve_zero, eq:impute,
%   eq:imputed_share, eq:cat_decomp, eq:effective_weight, eq:pillar, eq:psi.
% =============================================================================

\paragraph{Design of the analysis.} The pillars, the Strategic Index and the partnership
index are composites, and the uncertainty analysis asks how far the rankings they produce
depend on choices made in building them rather than on the data
\citep{saisana2005uncertainty, oecd2008handbook, saltelli2008global}. We vary one element at a
time, holding the others at the configuration of the final run: the vintage of the time-varying
sources, the category weights, the curvature $\gamma$ of the S-curve, the imputation rules, the
presence of each category, the rankability threshold $\tau$, the weights and terms of the
Strategic Index, the composition and eligibility rule of the Policy index, the rule that
screens Upstream picks for a resource base, the data choices made in the last revision of the
inputs, two candidate terms for vertical complementarities between stages, and the viewpoint
from which the partnership index is averaged. The comparison that matters for the pillars is the
ranking of countries within a technology--stage pair, so the main statistic is the Spearman rank
correlation between the baseline and the variant, computed within each pair and summarised
across pairs. Every comparison is made over the baseline's rankable cells
(Figure~\ref{fig:method_rankability}): $\ES$ is scored in 30 headline pairs, with 4{,}163
rankable cells of 6{,}232 scored, and $\EO$ in 33 pairs, with 4{,}301 rankable cells of 7{,}131.
Every variant is recomputed from the stored inputs rather than approximated from the published
indices. The weighting experiments use a re-implementation of the available-case mean of
equation~\eqref{eq:pillar} on the stored category scores, which reproduces the stored indices
at the configured weights to within $10^{-15}$; the curvature, imputation and data-choice
experiments rebuild both pillars from the theme tables with the builders and configuration of
the pipeline, which reproduce the stored indices exactly at the baseline settings; and the
Strategic Index and Policy index are rebuilt with the pipeline's own functions, which reproduce
the stored values exactly. Table~\ref{tab:robustness_summary} and
Figure~\ref{fig:robustness_variants} collect the results for the pillars.

% Generated by scripts/paper/fig_robustness.R; do not edit by hand.
\begin{table}[tbp]
  \centering
  \begin{threeparttable}
  \caption{Robustness of within-pair country rankings: Spearman correlations with the baseline}
  \label{tab:robustness_summary}
  \small
  \begin{tabular}{lcccccc}
    \toprule
     & \multicolumn{3}{c}{Energy Security ($\ES$)} & \multicolumn{3}{c}{Economic Opportunity ($\EO$)} \\
    \cmidrule(lr){2-4}\cmidrule(lr){5-7}
    Variant & Pooled & Median & Minimum & Pooled & Median & Minimum \\
    \midrule
    \multicolumn{7}{l}{\emph{Category weights}}\\
    \quad Dirichlet draws & 0.99 & 1.00 & 0.98 & 0.98 & 0.99 & 0.97 \\
    \quad Equal weights & 0.97 & 1.00 & 0.98 & 0.98 & 1.00 & 1.00 \\
    \multicolumn{7}{l}{\emph{S-curve curvature} (baseline $\gamma=0.5$)}\\
    \quad $\gamma=0.25$ & 0.98 & 1.00 & 0.97 & 0.98 & 0.99 & 0.98 \\
    \quad $\gamma=1$ (percent rank) & 0.98 & 0.99 & 0.90 & 0.98 & 0.99 & 0.97 \\
    \quad $\gamma=2$ & 0.93 & 0.96 & 0.74 & 0.94 & 0.97 & 0.90 \\
    \multicolumn{7}{l}{\emph{Global-average imputation rules}}\\
    \quad Removed (available case) & 0.99 & 1.00 & 0.62 & 0.85 & 0.97 & 0.29 \\
    \quad Replaced by zero & 0.87 & 0.98 & 0.58 & 0.62 & 0.77 & 0.42 \\
    \multicolumn{7}{l}{\emph{Vintage}}\\
    \quad 2020 against 2025 & 0.97 & 0.99 & 0.68 & 0.89 & 0.87 & 0.49 \\
    \quad 2025 vintage against headline & 0.99 & 1.00 & 0.87 & 0.99 & 1.00 & 0.79 \\
    \bottomrule
  \end{tabular}
  \begin{tablenotes}[flushleft]\footnotesize
    \item \emph{Notes:} Spearman rank correlations between the baseline pillar and each variant, over the baseline's rankable cells (sovereign economies with at most half of the index weight imputed). ``Pooled'' is computed over all rankable cells of the eleven headline technologies; ``Median'' and ``Minimum'' are taken across technology--stage pairs of the correlation computed within each pair. For the Dirichlet draws the pooled entry is the median over 1{,}000 draws and the within-pair entries summarise, across pairs, each pair's median over the draws. The vintage rows compare the 2020 and 2025 builds of \texttt{scripts/40\_build\_index\_vintages.R}, in which only the time-varying sources differ (4{,}139 $\ES$ and 4{,}301 $\EO$ rankable cells observed in both vintages), and the 2025 build with the headline index. Sample: 30 $\ES$ pairs (4{,}163 rankable cells) and 33 $\EO$ pairs (4{,}301 rankable cells). Entries are rounded to two decimals. Source: \texttt{scripts/paper/fig\_robustness.R}; \texttt{paper/results/robustness.json}.
  \end{tablenotes}
  \end{threeparttable}
\end{table}


\begin{figure}[tbp]
  \centering
  \includegraphics[width=\textwidth]{robustness_variants}
  \caption{Stability of within-pair country rankings under alternative constructions of the
  pillars. Each row is one variant; for each technology--stage pair we compute the Spearman
  rank correlation across countries between the baseline pillar and the variant, over the
  baseline's rankable cells (sovereign economies with at most half of the index weight
  imputed), and the row summarises these correlations across pairs: open circle, median; thick
  bar, interquartile range; thin line, minimum to maximum. For the Dirichlet weight draws the
  per-pair value is the median over 1{,}000 draws. ``Equal category weights'' gives every
  observed category the same weight. The S-curve rows rebuild both pillars from the theme tables
  with every component index re-curved to $\gamma=0.25$, 1 (the percent rank itself) or 2,
  against the baseline $\gamma=0.5$. ``No imputation'' removes every global-average rule of
  \texttt{config/missing\_data.yml}, so that uncovered countries drop the component by the
  available-case rule; ``Zero'' fills them with 0 instead. The last row compares the 2020 and
  2025 vintages of the pillars, in which only the time-varying sources differ. Sample: the
  eleven headline technologies; 30 $\ES$ pairs (4{,}163 rankable cells) and 33 $\EO$ pairs
  (4{,}301 rankable cells). Source: \texttt{index\_outputs.rds}, theme tables and
  \texttt{data/processed/vintages} of the final run; \texttt{scripts/paper/fig\_robustness.R}.}
  \label{fig:robustness_variants}
\end{figure}

\paragraph{Vintages.} The pillars are a cross-section in which each component takes the latest
year its source records (Section~\ref{subsec:norm_latest}). A separate build re-derives both
pillars as of 2020 and 2025 by slicing the time-varying sources (the EI Statistical Review, the
Atlas of Economic Complexity, GCIM investment, and the IMF and ILO price, rate and earnings
series) to each year and holding the themes without a time dimension (resource potential,
reserves, the IEA critical-mineral and clean-technology tables, the mineral production sheets,
BNEF demand, levelised costs and overcapacity, and technology readiness) identical. In this
build the Trade category is measured on the Atlas basis (HS92, annual) rather than on the
reconciled Comtrade pull of the headline. The comparison therefore measures how much the
rankings move with the information that can move; it is neither a measure of the total change
in countries' positions nor a validation. By the build's own classification, categories carrying
52 per cent of the nominal $\ES$ weight and 70 per cent of the $\EO$ weight are time-varying.
Not all of them change in practice: Consumption, in both pillars, and the $\EO$ Trade category
have identical scores in the two vintages, so the categories whose scores differ in at least one
per cent of their cells carry 40 per cent of the $\ES$ weight and 50 per cent of the $\EO$
weight. Because the weights are renormalised over the categories observed in each cell, the
share of a cell's weight that re-derives is lower still, 34 per cent in the average $\ES$ cell
and 36 per cent in the average $\EO$ cell, and 32 and 31 per cent once the Energy Prices
category is set aside, which is constant within every pair and so cannot reorder countries
(equation~\eqref{eq:affine_constant}).

Figure~\ref{fig:robustness_vintage} plots the 2025 against the 2020 value of each rankable
cell observed in both vintages (4{,}139 in $\ES$, 4{,}301 in $\EO$). $\ES$ moves little. The
median within-pair Spearman correlation is 0.99 (interquartile range 0.95 to 1.00), the mean
absolute change in the index is 0.015, 31 per cent of cells do not change at all, and a
country's percentile position in its pair moves by 3.8 points on average. Most of the visible
dispersion of the Upstream points is a parallel shift of whole pairs: 77 per cent of the sum of
squared Upstream changes is the common change of the pair, against 3 per cent Midstream and
4 per cent Downstream. Because the components are percent ranks, whose mean over a pair barely
moves, such a shift comes mainly from the category that is constant within the pair, Energy
Prices, whose world price-volatility series moves every economy of a pair by the same amount;
it changes levels and leaves ranks intact. $\EO$ moves more. The median within-pair
correlation is 0.87 (0.79 to 0.93), the mean absolute change is 0.030, no cell is unchanged,
and the mean percentile shift is 10.8 points. The difference between the pillars is a property
of the construction. Cost Competitiveness is scored in every $\EO$ pair and re-derives in all
its cells from the annual earnings, lending-rate and producer-price series, whereas the $\ES$
Upstream pairs are dominated by registers of reserves, production and mineral trade that the
vintage build holds fixed. In both pillars the lowest correlations are Downstream: Oil
Downstream is the minimum in each (0.68 in $\ES$, 0.49 in $\EO$), followed by Solar (0.69) and
Electric Vehicles (0.76) Downstream in $\ES$ and Gas (0.58) and Coal (0.69) Downstream in
$\EO$. The largest single moves occur where a pair rests on one flow variable. In Electric
Vehicles Downstream $\ES$ has a single category, Investment, built from annual GCIM flows, and
Viet Nam moves from 94th of 101 rankable economies to 3rd, and T\"urkiye from 39th to 1st. In
$\EO$ the largest coalition moves are Ukraine in Solar Downstream (42nd to 147th of 148) and Viet Nam in Electric Vehicles Downstream (41st to 2nd
of 105).

The 2025 vintage is close to, but not identical with, the headline index. Its category scores
coincide with the headline ones except in Trade, which the vintage build measures on the Atlas
basis, and its median within-pair correlation with the headline index is 1.00 in $\ES$ and
1.00 in $\EO$, with minima of 0.87 ($\ES$, Nuclear Midstream) and 0.79 ($\EO$, Nuclear
Upstream), two pairs whose trade baskets the final revision changed when it moved the uranium codes from the Midstream to the Upstream basket. Three limits
follow. Because only about a third of the average cell's weight re-derives, the vintage
comparison is likely to understate the instability of the rankings over time; because the
country lists of some sources differ slightly between years, part of what moves is coverage
rather than change; and the vintage years are nominal for two sources, whose nearest available
years are 2023 for the Atlas of Economic Complexity in the 2025 vintage and 2021 for the IMF
producer prices in the 2020 vintage.

\begin{figure}[tbp]
  \centering
  \includegraphics[width=\textwidth]{robustness_vintage}
  \caption{The 2020 and 2025 vintages of the pillars. Each point is a rankable country cell
  of the eleven headline technologies observed in both vintages (4{,}139 in $\ES$, 4{,}301 in
  $\EO$), coloured by supply-chain stage; the line is the 45-degree line. Both vintages are
  built by \texttt{scripts/40\_build\_index\_vintages.R} with the same theme builders, index
  builders and configuration as the headline run; only the time-varying sources (EI
  Statistical Review, Atlas of Economic Complexity, GCIM investment, IMF commodity prices,
  lending rates and producer prices, ILO earnings) are sliced to the vintage year, the Trade
  category is measured on the Atlas basis, and every other theme is held identical. Index
  values are on the unit interval; the indices are relative standings within each year's
  comparison sets, not absolute levels. Source:
  \texttt{data/processed/vintages/*\_comparison\_2020\_vs\_2025.csv};
  \texttt{scripts/paper/fig\_robustness.R}.}
  \label{fig:robustness_vintage}
\end{figure}

\paragraph{Category weights.} The weights of \texttt{config/weights.yml} are expert choices
\citep{decancq2013weights}. We draw 1{,}000 weight vectors from a Dirichlet distribution centred
on the configured shares, $w\sim\mathrm{Dir}(c\,\bar w)$, so that $E[w_k]=\bar w_k$ and the
coefficient of variation of $w_k$ is $\sqrt{(1-\bar w_k)/[\bar w_k(c+1)]}$. The concentration
$c=99$ for the ten categories of each pillar sets the coefficient of variation of a category of
average weight to 0.3; categories with a larger nominal weight vary less and smaller ones more.
For each draw both pillars are recomputed from the stored category scores with the
available-case rule. Within-pair rankings are almost insensitive to these perturbations. The
median over pairs of the within-pair Spearman correlation with the baseline has a median over
draws of 0.997 in $\ES$ (5th to 95th percentile over draws, 0.992 to 0.999) and 0.988 in
$\EO$ (0.960 to 0.997). Equal category weights give median within-pair correlations of
1.000 in both pillars, because the configured weights are themselves close to equal.

The insensitivity is largely a property of the construction rather than evidence that the
categories agree. Within every pair all scored countries observe the same category set
(Section~\ref{subsec:norm_pillar}), so by equation~\eqref{eq:affine_constant} a category that
is constant across them moves every country's index by the same affine map, and the ranking
depends only on the relative weights of the categories that vary. The median pair has three
varying categories in $\ES$ and four in $\EO$. In seven $\ES$ pairs and three $\EO$ pairs a
single category varies, and there the ranking does not depend on the weights at all; on
average 3 per cent of the observed weight of an $\ES$ pair and 25 per cent of an $\EO$ pair
sits on categories that are constant within it, among them Energy Prices and, in $\EO$,
Technological Readiness. Restricted to the pairs with two or more varying categories (23 in
$\ES$, 30 in $\EO$), the median correlation over pair--draws is still 0.996 in $\ES$ and 0.987
in $\EO$, but the minimum over pair--draws falls to 0.77 and 0.30.

A high average correlation can hide large moves of individual countries, which matter for
statements about coalition members. For each of the 57 members and each pair in which it is
rankable we record the 5th and 95th percentiles of its rank over the draws. In percentile
points of the pair, the median width of this 90 per cent band is 2.8 in $\ES$ (3 ranks, over
1{,}606 member cells) and 9.6 in $\EO$ (11 ranks, over 1{,}718); 79 per cent of member cells in
$\ES$ and 52 per cent in $\EO$ have a band narrower than ten points, but the 90th percentile of
the width is 18 points in $\ES$ and 31 in $\EO$. Wide bands arise where a member sits among
large blocks of tied scores. In $\ES$ Batteries Downstream, South Korea ranks 69th of 186 at the
configured weights, and its 90 per cent band runs from 15th to 186th, the widest in the sample.
In $\EO$ the widest bands are Upstream: Brazil in Nuclear Upstream (163rd of 175, band 21st to
174th) and in the oil and battery-mineral pairs, and Germany in Nuclear Upstream.
Figure~\ref{fig:robustness_ranks} shows the bands of the core members for one pair.

\begin{figure}[tbp]
  \centering
  \includegraphics[width=\textwidth]{robustness_ranks}
  \caption{Rank uncertainty of the core coalition members under perturbed category weights,
  Batteries Midstream. For each core member (ISO3 code) with a rankable cell in the pair (31 in
  $\ES$, 28 in $\EO$) the dot is its percentile position among the pair's rankable cells at the
  configured weights (100 = highest index) and the bar spans the 5th to 95th percentiles of that
  position over 1{,}000 Dirichlet draws of the category weights centred on the configured
  shares (coefficient of variation 0.3 for a category of average weight). Members are ordered by
  their baseline $\EO$ position; members without a rankable $\EO$ cell are placed at the
  bottom. A short bar means the member's standing does not depend on the weights within that
  range. Source: category scores in \texttt{index\_outputs.rds};
  \texttt{scripts/paper/fig\_robustness.R}.}
  \label{fig:robustness_ranks}
\end{figure}

\paragraph{Curvature of the S-curve.} Most component indices are an S-curve of a percent
rank, $\phi_{0.5}(r)$, and by equation~\eqref{eq:scurve_logit} $\operatorname{logit}\phi_\gamma(r)=
\gamma\operatorname{logit}r$. Such a component at another curvature $\gamma'$ can therefore be
obtained from the stored value without the raw data,
$\phi_{\gamma'}(r)=\operatorname{logistic}\bigl((\gamma'/0.5)\operatorname{logit}\phi_{0.5}(r)
\bigr)$, a map that also sends the anchored zeros of equation~\eqref{eq:scurve_zero} to zero.
We apply it to every index row of every theme table and rebuild both pillars for
$\gamma'\in\{0.25,1,2\}$. The value $\gamma'=1$ is the plain percent rank, and $\gamma'=2$ is
an S-shape that stretches the middle of the distribution and compresses the tails, the opposite
of the baseline \citep{karmarkar1978subjectively}. The map is exact for a component that is a
single S-curve of a rank. For components that a theme forms as a mean of several S-curves, or
by another transform, it re-curves the mean rather than its parts and is an approximation. We
classify each component by inverting $\phi_{0.5}$ and testing whether the implied ranks lie on a
rank grid. Of the 57 components that enter a scored category, 46 pass; they carry on average
84 per cent of the weight of an $\ES$ cell and 63 per cent of an $\EO$ cell. The exceptions are
the input-cost and levelised-cost terms of Cost Competitiveness, the technology-readiness and
overcapacity terms, the production surplus or deficit of Energy Imports, the mineral-supply and
overall-production shares of Foreign Dependency, the two minerals-trade terms, and the
geothermal and wind potential indices.

Within pairs the rankings are stable under the stronger inverse-S curve and under the plain
percent rank, and less so under the S-shape. The median within-pair correlation with the
baseline is 0.997 in $\ES$ and 0.992 in $\EO$ at $\gamma'=0.25$, 0.989 and 0.991 at
$\gamma'=1$, and 0.964 and 0.965 at $\gamma'=2$, where the minima are 0.74 ($\ES$, Batteries
Downstream) and 0.90 ($\EO$, Electric Vehicles Midstream). Because the map is monotone, a pair
whose ranking rests on a single varying component is unaffected; the curvature matters only
through the implicit weighting among components, a larger $\gamma$ giving more weight to
differences near the median and less to differences in the tails. Levels move more than ranks:
the mean absolute change of $\ES$ is 0.03, 0.04 and 0.08 for the three curvatures.

\paragraph{Imputation rules.} Nine $\ES$ themes and twelve $\EO$ themes fill uncovered countries
with the observed mean of their component group (equation~\eqref{eq:impute}). In the baseline,
components filled in this way carry on average 7 per cent of the weight of an $\ES$ cell and
25 per cent of an $\EO$ cell. In $\EO$ two thirds of that weight is in Technology Demand and
most of the rest in Production and Cost Competitiveness (16 and 10 per cent), categories whose
country-specific sources cover only part of the economies in each pair; in $\ES$ it is mainly in
Energy Access and Foreign Dependency (70 and 17 per cent). We rebuild both pillars with every
global-average rule removed, so that an uncovered country drops the component under the
available-case rule, and with every such rule replaced by zero. Over the rankable cells $\ES$ is
robust to the rule in most pairs: the median within-pair correlation is 1.00 without imputation
and 0.98 with zeros, with minima of 0.62 (Nuclear Upstream) and 0.58 (Coal Downstream). $\EO$ is
not robust to it. Without imputation its median within-pair correlation is 0.97 but the pooled
correlation 0.85 and the minimum 0.29 (Coal Midstream); with zeros the median falls to 0.77,
with a pooled correlation of 0.62. Zero-filling a narrow source places every uncovered economy
at or below every covered one on that component, so that within a pair the variant ranks
countries partly by whether a source covers them. Removing the rule scores an uncovered economy
on the components it does have, so that its index rests on a different and smaller set of
measurements than its neighbours'. The global-average rule places it at the mean of the covered
economies on the missing component, which compresses the differences that component would
make. Neither is neutral, and the configuration chooses between them theme by theme
(Section~\ref{subsec:norm_missing}). What the experiment shows is that the $\EO$ rankings,
outside the economies the narrow sources cover, depend on that choice; this is the reason the
rankability rule withholds from cross-country rankings the cells whose index rests mostly on
imputed components (equation~\eqref{eq:imputed_share}).

\paragraph{The rankability threshold.} A cell enters cross-country rankings, top-$N$ lists and
the Strategic Index only if the economy is sovereign and at most a share $\tau=0.5$ of its index
weight is imputed. Table~\ref{tab:robustness_rankability} varies $\tau$. Because the rule
filters cells and does not change values, the order of any two cells that are rankable under
two thresholds is the same, and what $\tau$ changes is the sample. The sample is sensitive to
it, $\EO$ more than $\ES$: at $\tau=0.3$ the rankable cells fall to 3{,}147 in $\ES$ (56 per
cent of the sovereign scored cells, against 74 per cent at the headline threshold) and to
1{,}967 in $\EO$ (31 against 67 per cent), and at $\tau=0.7$ they rise to 5{,}350 and 5{,}974
(95 and 93 per cent). The threshold also governs how far the rankings depend on the imputation
rule. At $\tau=0.3$ the median within-pair correlation of $\EO$ with its zero-filled variant is
0.94; at the headline threshold it is 0.77 and at $\tau=0.7$ 0.68. The headline lists are
stable to a stricter threshold in $\ES$ (the top ten of 27 of the 30 pairs are identical at
$\tau=0.3$) but not in $\EO$, where 15 of the 33 top-ten lists are identical and the mean
overlap is 77 per cent, because the stricter threshold removes economies whose $\EO$ rests on
regional demand outlooks and cost terms imputed from group means. The Strategic Index, which
needs rankable $\ES$ and $\EO$, shrinks from 3{,}154 cells for 174 countries to 1{,}577 cells for
107 countries at $\tau=0.3$, where 48 of the 57 coalition members keep a pick and 29 of them the
baseline one; at $\tau=0.7$ it has 5{,}156 cells for 194 countries and 42 of the 57 members
keep the baseline pick. The threshold is thus a substantive choice for the allied selection, which
we report with the headline value and these alternatives rather than treat as settled.

% Generated by scripts/paper/fig_robustness.R; do not edit by hand.
\begin{table}[tbp]
  \centering
  \begin{threeparttable}
  \caption{The rankability threshold: cells ranked, robustness of their rankings, and the Strategic Index}
  \label{tab:robustness_rankability}
  \small
  \begin{tabular}{lcccccccccc}
    \toprule
    & \multicolumn{3}{c}{$\ES$} & \multicolumn{3}{c}{$\EO$} & \multicolumn{3}{c}{Strategic Index} & PSI \\
    \cmidrule(lr){2-4}\cmidrule(lr){5-7}\cmidrule(lr){8-10}\cmidrule(lr){11-11}
    $\tau$ & Cells & $\rho_{\text{none}}$ & $\rho_{\text{zero}}$ & Cells & $\rho_{\text{none}}$ & $\rho_{\text{zero}}$ & Cells & Countries & Picks & Ranked \\
    \midrule
    0.3 & 3{,}147 & 1.00 & 1.00 & 1{,}967 & 0.99 & 0.94 & 1{,}577 & 107 & 29/48 & 159/102 \\
    0.5$^{\ast}$ & 4{,}163 & 1.00 & 0.98 & 4{,}301 & 0.97 & 0.77 & 3{,}154 & 174 & 57/57 & 192/188 \\
    0.7 & 5{,}350 & 0.99 & 0.98 & 5{,}974 & 0.86 & 0.68 & 5{,}156 & 194 & 42/57 & 194/194 \\
    \bottomrule
  \end{tabular}
  \begin{tablenotes}[flushleft]\footnotesize
    \item \emph{Notes:} A cell is rankable when the economy is sovereign and at most a share $\tau$ of its index weight rests on imputed components ($^{\ast}$ the headline, \texttt{config/rankability.yml}). Cells: rankable headline cells (of 6{,}232 scored $\ES$ and 7{,}131 scored $\EO$ cells). $\rho_{\text{none}}$ and $\rho_{\text{zero}}$: median across technology--stage pairs of the within-pair Spearman correlation, over the cells rankable at $\tau$, between the baseline pillar and the pillar rebuilt with every global-average rule removed (available case) or replaced by zero. Strategic Index: cells and countries in the sample when the pillars' rankable flags are taken at $\tau$ (the Policy index is unchanged), and the number of coalition members whose allied pick is the baseline one, over the members with a pick at $\tau$. PSI: partners ranked in the core-coalition (headline) and world views. Within a technology--stage pair the order of two cells rankable under two thresholds is the same, because the rule filters cells and does not change values. Source: \texttt{scripts/paper/fig\_robustness.R} (\texttt{rankability} in \texttt{paper/results/robustness.json}).
  \end{tablenotes}
  \end{threeparttable}
\end{table}


\paragraph{Leaving out one category.} Setting one category's weight to zero and renormalising
the others measures how much of the within-pair ranking each category carries
(Figure~\ref{fig:robustness_loo}). We report the correlation only over the pairs in which the
category is observed, since elsewhere dropping it changes nothing. In $\ES$ the ranking depends
most on Consumption (median within-pair correlation 0.74 over the eight pairs in which it is
observed alongside other categories; minimum 0.46 in Batteries Downstream) and on Minerals Trade
(0.79 over six pairs). Energy Access, Energy Imports, Investment, Trade and Reserves follow, with
medians between 0.89 and 0.94; Trade has the widest range (minimum 0.40 in Electric Grid
Midstream, whose basket the completed trade file filled), and in the four Midstream fossil and nuclear pairs it is the
only category, so that dropping it leaves the pair unscored. Dropping Production or Foreign
Dependency, which share the largest nominal $\ES$ weight, changes almost nothing (medians 0.995
and 0.998): Production, observed in seven pairs, separates a few large producers from the rest,
and a rank correlation over some two hundred economies barely registers the reordering of a few
leaders. Energy Prices, constant within every pair, changes nothing at
all. In $\EO$ three categories carry most of the ranking: Foreign Dependency (median 0.75 over
its five pairs), Cost Competitiveness (0.76 over the 30 pairs in which other categories vary)
and Trade (0.79 over 18 pairs). Technological Readiness and Energy Prices are cell constants and
leave every ranking unchanged, and Technology Demand, the category with the largest nominal
$\EO$ weight, changes the median ranking almost not at all (0.99). As in other composite
indicators, nominal weights and effective importance diverge
\citep{paruolo2013ratings, becker2017weights}: the categories with the largest nominal weights
are not the ones that rank countries within a pair.

\begin{figure}[tbp]
  \centering
  \includegraphics[width=\textwidth]{robustness_loo}
  \caption{Leave-one-category-out sensitivity of within-pair rankings. Each row drops one
  category (weight set to zero, the remaining observed categories renormalised by the
  available-case rule) and reports the Spearman rank correlation across countries between the
  baseline pillar and the reduced pillar, computed within each technology--stage pair in which
  the category is observed, over the baseline's rankable cells: open circle, median over those
  pairs; thick bar, interquartile range; thin line, minimum to maximum. The label gives the
  category's nominal weight share and the number of headline pairs in which it is observed.
  Pairs in which the dropped category is the only varying one are omitted, because the reduced
  pillar is then constant or unscored (four $\ES$ pairs for Trade, three $\EO$ pairs for Cost
  Competitiveness, one $\ES$ pair each for Consumption, Foreign Dependency and Investment). Rows
  are ordered by the median. Source: category scores in \texttt{index\_outputs.rds};
  \texttt{scripts/paper/fig\_robustness.R}.}
  \label{fig:robustness_loo}
\end{figure}

\paragraph{Nominal and effective importance.} The leave-one-out correlations measure how much
of the ranking a category carries. A complementary measure, a linear and exact counterpart of
the main effects that \citet{paruolo2013ratings} propose as the effective importance of an
indicator, is the share of the cross-country variance of a pillar within a pair for which each
category accounts. Because the pillar is the sum of its category contributions
$C_{kits}=\omega_{kits}c_{kits}$ (equation~\eqref{eq:cat_decomp}), its variance across the
rankable countries of a pair splits exactly into covariance shares,
\begin{equation}\label{eq:variance_share}
  S_{kts}=\frac{\operatorname{Cov}_i\bigl(C_{kits},\,\ES_{its}\bigr)}
               {\operatorname{Var}_i\bigl(\ES_{its}\bigr)},
  \qquad \sum_{k\in\Cats^{P}}S_{kts}=1,
\end{equation}
and likewise for $\EO$. A category that is constant within the pair has $S_{kts}=0$ whatever
its weight. Table~\ref{tab:robustness_importance} compares, for each category, the nominal
weight share, the effective weight $\omega_{kits}$ of equation~\eqref{eq:effective_weight} and
the variance share, the last two averaged over the headline pairs with zero where the category
is not observed, so that each column sums to one within a pillar.

The three measures differ considerably. Half the sum of the absolute differences between the
nominal shares and the mean variance shares, the share of weight that would have to be moved
for the two to coincide, is 0.32 in $\ES$ and 0.37 in $\EO$, but the gap has different sources
in the two pillars. In $\ES$ it is largely a matter of coverage, of which categories are
observed in which pairs: the corresponding distance between the effective weights and the
variance shares is 0.19. Trade, the only category of the four Midstream fossil and nuclear
pairs and now observed in 18 of the 30 pairs (the new Upstream baskets added four), has
9.5 per cent of the nominal weight but 24.9 per cent of the mean effective weight and 21.5 per
cent of the variance, while Production and Energy Imports, observed only in seven and three
pairs, fall from nominal shares of 14.3 and 9.5 per cent to variance shares of 10.0 and
2.8 per cent. In $\EO$ the gap arises within pairs, and the distance between effective weights
and variance shares is 0.38. Technology Demand and Technological Readiness hold 21.8 and 14.6
per cent of the mean effective weight but 3.6 and 0.0 per cent of the variance, because the
readiness score is a technology-level constant and the only country-varying components of
Technology Demand are the regional demand outlooks, whereas Cost Competitiveness, observed in
every pair, holds 21.9 per cent of the effective weight and 32.8 per cent of the variance. Which
categories matter therefore depends on the question: the variance shares describe what
separates the leaders of a pair from the rest, the leave-one-out correlations what orders the
bulk of the distribution.

% Generated by scripts/paper/fig_robustness.R; do not edit by hand.
\begin{table}[tbp]
  \centering
  \begin{threeparttable}
  \caption{Nominal weight, effective weight and variance share of each category within technology--stage pairs}
  \label{tab:robustness_importance}
  \small
  \begin{tabular}{lrrrrr}
    \toprule
    & Nominal & Effective & Variance & & Leave-one- \\
    Category & share (\%) & weight (\%) & share (\%) & Pairs & out $\rho$ \\
    \midrule
    \multicolumn{6}{l}{\emph{Panel A. Energy Security ($\ES$), 30 pairs}}\\
    \quad Trade & 9.5 & 24.9 & 21.5 & 18 & 0.93 \\
    \quad Investment & 7.1 & 11.3 & 18.1 & 12 & 0.92 \\
    \quad Reserves & 9.5 & 9.5 & 17.4 & 13 & 0.94 \\
    \quad Consumption & 11.9 & 15.8 & 13.2 & 9 & 0.74 \\
    \quad Foreign Dependency & 14.3 & 16.1 & 10.1 & 12 & 1.00 \\
    \quad Production & 14.3 & 5.8 & 10.0 & 7 & 0.99 \\
    \quad Energy Access & 9.5 & 7.2 & 5.0 & 6 & 0.89 \\
    \quad Energy Imports & 9.5 & 2.9 & 2.8 & 3 & 0.91 \\
    \quad Minerals Trade & 9.5 & 3.1 & 1.8 & 6 & 0.79 \\
    \quad Energy Prices & 4.8 & 3.4 & 0.0 & 10 & 1.00 \\
    \addlinespace
    \multicolumn{6}{l}{\emph{Panel B. Economic Opportunity ($\EO$), 33 pairs}}\\
    \quad Cost Competitiveness & 10.0 & 21.9 & 32.8 & 33 & 0.76 \\
    \quad Trade & 10.0 & 11.0 & 17.5 & 18 & 0.79 \\
    \quad Investment & 10.0 & 6.6 & 14.4 & 12 & 0.90 \\
    \quad Production & 10.0 & 8.1 & 12.5 & 16 & 0.95 \\
    \quad Consumption & 10.0 & 5.0 & 7.9 & 9 & 0.91 \\
    \quad Foreign Dependency & 10.0 & 2.9 & 7.4 & 5 & 0.75 \\
    \quad Energy Access & 10.0 & 2.8 & 3.8 & 6 & 0.91 \\
    \quad Technology Demand & 13.3 & 21.8 & 3.6 & 27 & 0.99 \\
    \quad Energy Prices & 10.0 & 5.3 & 0.0 & 10 & 1.00 \\
    \quad Technological Readiness & 6.7 & 14.6 & 0.0 & 33 & 1.00 \\
    \bottomrule
  \end{tabular}
  \begin{tablenotes}[flushleft]\footnotesize
    \item \emph{Notes:} Nominal share: $w_k/\sum_{k'}w_{k'}$ from \texttt{config/weights.yml}. Effective weight: mean over the rankable headline pairs of the weight $w_k/\sum_{l\ \mathrm{observed}}w_l$ a category receives after renormalisation over the observed categories; zero where the category is not observed. Variance share: mean over the same pairs of $S_{kts}$ (equation~\eqref{eq:variance_share}), the covariance of the category's contribution with the pillar across the rankable cells of the pair divided by the pillar's variance; zero where the category is not observed. Each of the three columns sums to 100 within a panel up to rounding. Pairs: number of pairs in which the category is observed. Leave-one-out $\rho$: median, over the pairs in which the category is observed and the reduced pillar still ranks countries, of the within-pair Spearman correlation between the baseline pillar and the pillar without the category (Figure~\ref{fig:robustness_loo}). Rows are ordered by the variance share. Sample: rankable cells of the eleven headline technologies. Source: category scores in \texttt{index\_outputs.rds}; \texttt{scripts/paper/fig\_robustness.R} (\texttt{weights.<pillar>.variance\_importance} in \texttt{paper/results/robustness.json}).
  \end{tablenotes}
  \end{threeparttable}
\end{table}


\paragraph{Strategic Index: weights and specification.} The same Dirichlet design applied to
the five weights of equation~\eqref{eq:si} (concentration 43.4, so that the coefficient of
variation is 0.30 for the policy weight, 0.23 for the two pillar weights and 0.45 for the two
sector priors) leaves the ordering of each country's sectors close to the baseline: over the
3{,}154 cells of the 174 countries with a Strategic Index, the median over countries of the
within-country Spearman correlation has a median of 0.96 over draws (5th to 95th percentile
0.88 to 0.99). The allied pick, a single arg max per member over its eligible Upstream and
Midstream sectors (equation~\eqref{eq:si_allied}), is more fragile
(Figure~\ref{fig:robustness_si}, panel a). On average a member keeps its baseline pick in 73
per cent of the draws (median 82 per cent); 16 of the 57 members keep it in at least 90 per cent
and 10 in fewer than half, the least stable being Indonesia (15 per cent), whose two leading
sectors differ by 0.0007. Retention is close to a monotone function of the baseline margin
between the first and second eligible sectors (Spearman correlation 0.84 across members). The
margins are small: their median is 0.015 (interquartile range 0.006 to 0.024), and the margin is
below 0.005 for 10 of the 57 members, below 0.01 for 21 and below 0.02 for 37. They differ by
sector. The four Electric Grid Midstream picks have a median margin of 0.003 and the nine
Nuclear Midstream picks 0.007, against 0.014 for the 27 Batteries Midstream picks and 0.024 for
the eight Nuclear Upstream picks, so the picks that most depend on fine margins are the grid and
Midstream nuclear ones.

The members with the largest margins are not those with the clearest comparative advantage. Costa
Rica, the Dominican Republic and Guatemala, whose Nuclear Midstream picks keep their place in
nearly every draw, have only four eligible sectors each (the Midstream stages of nuclear, oil,
gas and coal), and among these the technology greenhouse-gas prior favours nuclear; without the
two sector priors the Dominican Republic and Guatemala move to Oil Midstream. Their picks are
therefore a consequence of a narrow rankable sample and of the prior, not of an $\ES$ floor or of
a measured nuclear capability, and we read them as such.

Panel (b) of Figure~\ref{fig:robustness_si} and Panel A of Table~\ref{tab:robustness_picks}
report the discrete variants (see also Figure~\ref{fig:strategic_sign}). Reading security as
strength rather than need is the most consequential choice: the pooled Spearman correlation with
the baseline is 0.33, the median within-country correlation 0.32, and the allied pick survives
for only 5 of the 57 members; 28 members then pick Nuclear Midstream and 10 Gas Midstream.
The other specification variants preserve the ordering of sectors much better, with median
within-country correlations of 0.86 without the sector priors, 0.85 with $\EO$ and $1-\ES$
S-curved within country like the policy term, 0.85 without the policy term and 0.92 with equal
weights on the five terms, but keep only 31, 20, 30 and 39 of the 57 picks. The direction of the
changes is informative. Without the priors, ten picks move to Oil Midstream, the sector that the
technology prior penalises most; without the policy term, 20 members pick Nuclear Upstream
instead of 8, so that for many members with recorded uranium resources the pillars alone point
to uranium extraction and the policy term pulls the pick elsewhere. The ordering of a member's sectors is thus robust to moderate changes in the
weights, while the identity of its single top sector, which typically turns on margins of one or
two hundredths, is robust only where that margin is large.

\begin{figure}[tbp]
  \centering
  \includegraphics[width=\textwidth]{robustness_si}
  \caption{Robustness of the allied picks. Each row is a coalition member (ISO3 code), its
  baseline pick, the Upstream or Midstream sector (excluding Geothermal and Green Hydrogen, and
  Upstream sectors without a recorded resource base) with the highest Strategic Index
  (equation~\eqref{eq:si_allied}), and in brackets the margin of its Strategic Index over the
  member's second eligible sector. Panel (a): share of 1{,}000 Dirichlet draws of the five
  Strategic Index weights, centred on $(0.2,0.3,0.3,0.1,0.1)$ with concentration 43.4, under
  which the member's pick is unchanged; colour gives the stage of the baseline pick. Panel (b):
  filled circle, the pick is unchanged under the variant; open circle, it changes; cross, the
  member has no pick under the variant; the count at the foot of each column is the number of
  unchanged picks. Variants: security $\ES_{its}$ in place of need $1-\ES_{its}$; without the
  stage and technology priors; $\EO_{its}$ and $1-\ES_{its}$ replaced by their within-country
  S-curves; without the policy term; both Policy components required
  (\texttt{OPSI\_POLICY\_REQUIRE\_ALL=1}), under which the 11 members scored on IEA PAMS alone
  have no Policy index and so no pick; CAT dropped for every country (IEA PAMS only); the October
  2026 NIPO component added back to the Policy index with weight 4 (IEA PAMS 3, CAT 2); the
  endowment screen switched off; and the Trade category of both pillars measured on the Atlas
  basis of the 2025 vintage. Sample: 3{,}154 Strategic Index cells of 174 countries; 57 coalition
  members, all with a baseline pick. Source: \texttt{strategic\_index}, \texttt{policy\_component\_tbl}
  and category scores in \texttt{index\_outputs.rds}; \texttt{<scratch>/nipo\_oct};
  \texttt{scripts/paper/fig\_robustness.R}.}
  \label{fig:robustness_si}
\end{figure}

\paragraph{The endowment screen.} An Upstream extraction stage is eligible for a member's pick
only if the country has recorded reserves or production of that stage's resources; the screen
leaves the Strategic Index itself unchanged and marks 669 of the 998 scored Upstream cells as
having no resource base. Switching it off changes 41 of the 57 picks: the Upstream picks rise
from 14 to 47, Nuclear Upstream from 8 to 17 and Solar Upstream from none to 12, and members
with no uranium, among them Luxembourg, the Netherlands, Ireland and Denmark, would pick Nuclear
Upstream. The reason is a floor in the security pillar. In several sectors a large share of
countries share one $\ES$ value, because the categories that vary are observed only for
producers: 34 of the 40 countries in Green Hydrogen Midstream, 86 of 174 in Nuclear Upstream,
45 of 158 in Green Hydrogen Downstream, 19 of 67 in Nuclear Downstream and 38 of 138 in
Batteries Downstream sit at their sector's modal value. In Nuclear Upstream that common value,
0.17, is close to the pair's minimum of 0.14: a country without uranium scores zero on the
reserve and production terms, so it receives a low $\ES$ and a high security need $1-\ES$, and with the
nuclear technology prior this makes Nuclear Upstream the top Upstream sector of many countries
that could not supply it. The screen is a rule of eligibility for the pick, not part of the index,
and the comparison shows that without it the Upstream picks would largely reflect that floor.

% Generated by scripts/paper/fig_robustness.R; do not edit by hand.
\begin{table}[tbp]
  \centering
  \begin{threeparttable}
  \caption{Allied picks by sector under the baseline and the variants of the Strategic Index}
  \label{tab:robustness_picks}
  \footnotesize
  \setlength{\tabcolsep}{4pt}
  \begin{tabular}{lrrrrrrrrr}
    \toprule
    \multicolumn{10}{l}{\emph{Panel A. Specification and Policy index}}\\
    Sector & \makecell{Base-\\line} & \makecell{$\ES$ for\\$1-\ES$} & \makecell{No\\priors} & \makecell{All\\ranked} & \makecell{No\\policy} & \makecell{Equal\\weights} & \makecell{CAT\\required} & \makecell{PAMS\\only} & \makecell{NIPO\\added} \\
    \midrule
    Nuclear Upstream & 8 & 2 & 14 & 13 & 20 & 1 & 6 & 4 & 14 \\
    Batteries Upstream & 4 & 1 & 3 & 4 & 2 & 2 & 3 & 4 & 6 \\
    Wind Upstream & 2 & 0 & 2 & 0 & 3 & 1 & 1 & 1 & 3 \\
    Coal Upstream & 0 & 0 & 1 & 0 & 0 & 0 & 0 & 0 & 0 \\
    Electric Grid Upstream & 0 & 3 & 1 & 1 & 1 & 0 & 0 & 0 & 0 \\
    Electric Vehicles Upstream & 0 & 0 & 1 & 0 & 0 & 0 & 0 & 0 & 0 \\
    Gas Upstream & 0 & 1 & 0 & 0 & 0 & 0 & 0 & 0 & 0 \\
    Oil Upstream & 0 & 2 & 3 & 0 & 0 & 0 & 0 & 0 & 0 \\
    Solar Upstream & 0 & 0 & 2 & 4 & 1 & 0 & 0 & 0 & 2 \\
    Batteries Midstream & 27 & 0 & 15 & 5 & 8 & 30 & 26 & 34 & 15 \\
    Nuclear Midstream & 9 & 28 & 2 & 8 & 9 & 8 & 5 & 5 & 7 \\
    Electric Grid Midstream & 4 & 7 & 1 & 15 & 13 & 4 & 3 & 3 & 5 \\
    Solar Midstream & 3 & 2 & 2 & 1 & 0 & 9 & 2 & 6 & 4 \\
    Electric Vehicles Midstream & 0 & 0 & 0 & 0 & 0 & 0 & 0 & 0 & 1 \\
    Gas Midstream & 0 & 10 & 0 & 0 & 0 & 0 & 0 & 0 & 0 \\
    Oil Midstream & 0 & 1 & 10 & 6 & 0 & 0 & 0 & 0 & 0 \\
    Wind Midstream & 0 & 0 & 0 & 0 & 0 & 2 & 0 & 0 & 0 \\
    \midrule
    Members with a pick & 57 & 57 & 57 & 57 & 57 & 57 & 46 & 57 & 57 \\
    Same pick as baseline & 57 & 5 & 31 & 20 & 30 & 39 & 46 & 45 & 34 \\
    \addlinespace
    \midrule
    \multicolumn{10}{l}{\emph{Panel B. Data: endowment screen, final data choices, trade and reserves inputs, vintage}}\\
    Sector & \makecell{Base-\\line} & \makecell{Screen\\off} & \makecell{Zero for\\EIA} & \makecell{No new\\baskets} & \makecell{Trade on\\Atlas} & \makecell{No\\Trade} & \makecell{No\\Reserves} & \makecell{2025\\vintage} & \makecell{2020\\vintage} \\
    \midrule
    Nuclear Upstream & 8 & 17 & 8 & 11 & 8 & 9 & 14 & 9 & 8 \\
    Batteries Upstream & 4 & 12 & 4 & 4 & 4 & 1 & 1 & 4 & 5 \\
    Wind Upstream & 2 & 5 & 2 & 1 & 2 & 1 & 2 & 2 & 1 \\
    Electric Grid Upstream & 0 & 1 & 0 & 0 & 0 & 0 & 0 & 0 & 0 \\
    Solar Upstream & 0 & 12 & 0 & 0 & 0 & 0 & 0 & 0 & 0 \\
    Batteries Midstream & 27 & 7 & 27 & 25 & 26 & 13 & 24 & 25 & 23 \\
    Nuclear Midstream & 9 & 2 & 9 & 9 & 4 & 0 & 8 & 4 & 4 \\
    Electric Grid Midstream & 4 & 1 & 4 & 4 & 5 & 10 & 4 & 5 & 11 \\
    Solar Midstream & 3 & 0 & 3 & 3 & 7 & 19 & 4 & 7 & 5 \\
    Electric Vehicles Midstream & 0 & 0 & 0 & 0 & 0 & 1 & 0 & 0 & 0 \\
    Gas Midstream & 0 & 0 & 0 & 0 & 1 & 0 & 0 & 1 & 0 \\
    \midrule
    Members with a pick & 57 & 57 & 57 & 57 & 57 & 54 & 57 & 57 & 57 \\
    Same pick as baseline & 57 & 16 & 57 & 53 & 49 & 29 & 50 & 48 & 44 \\
    \bottomrule
  \end{tabular}
  \begin{tablenotes}[flushleft]\scriptsize
    \item \emph{Notes:} Number of coalition members whose allied pick, the Upstream or Midstream sector (excluding Geothermal and Green Hydrogen, and Upstream cells without a recorded resource base) with the highest Strategic Index, falls in each sector. All 57 members have a pick in the baseline. Panel A: $\ES$ for $1-\ES$, security in place of need; No priors, without the stage and technology terms; All ranked, $\EO$ and $1-\ES$ S-curved within country; No policy, without the policy term; Equal weights, 0.2 on each term; CAT required, both Policy components required (\texttt{OPSI\_POLICY\_REQUIRE\_ALL=1}), which leaves 65 countries with a Policy index and 46 members with a pick; PAMS only, CAT dropped for every country; NIPO added, the October 2026 NIPO component (GTA NIPO as of 5 October 2026) re-enters the Policy index with weight 4 (IEA PAMS 3, CAT 2), IEA PAMS still required, for its 77 policy-eligible jurisdictions. Panel B: Screen off, every Upstream cell eligible, including the 669 of 998 scored Upstream cells without recorded reserves or production; Zero for EIA, the oil, gas and coal production that the Energy Institute does not report set to zero instead of the EIA value (the rule until 11 October 2026); No new baskets, the Trade category removed from Electric Vehicles, Wind, Nuclear and Solar Upstream, whose trade baskets were added on 11 October 2026; Trade on Atlas, the Trade category of both pillars replaced by its Atlas-based counterpart of the 2025 vintage; No Trade, No Reserves, the category dropped from the pillars that score it and the remaining categories renormalised; 2025 and 2020 vintage, both pillars taken from the vintage builds, in which only the time-varying sources differ. The Policy index is the headline one in Panel B. Source: \texttt{index\_outputs.rds}, \texttt{data/processed/vintages}; \texttt{scripts/paper/fig\_robustness.R}.
  \end{tablenotes}
  \end{threeparttable}
\end{table}


\paragraph{The Policy index.} Since 9 October 2026 the Policy index requires IEA PAMS and adds
CAT where it rates the country, so that 215 countries are scored, 65 on both components and
150 on IEA PAMS alone, among them 11 of the 57 coalition members. Three variants test this
choice (Table~\ref{tab:robustness_data} and Panel A of Table~\ref{tab:robustness_picks}). First,
requiring both components (\texttt{OPSI\_POLICY\_REQUIRE\_ALL=1}, the earlier rule) leaves 65
countries with a Policy index and a Strategic Index of 1{,}666 cells. Because the composite is
renormalised over the components present in each cell, a country with both components has the
same Policy index under either rule, and its Strategic Index is identical; the 46 members with a
CAT rating therefore keep their picks exactly, and the 11 members scored on IEA PAMS alone lose
theirs. The rule decides who has a pick, not what the pick is. Second, the basis matters for
the values: for the 65 countries with both components, the within-country Spearman correlation
between the policy values on IEA PAMS alone and on the composite has a median of 0.48
(interquartile range 0.41 to 0.57). Dropping CAT for every country nevertheless keeps the
Strategic Index ordering largely intact (median within-country correlation 1.00, interquartile
range 0.91 to 1.00, minimum 0.30) and 45 of the 57 picks, because policy carries a fifth of
the index. Third, the NIPO component of the Global Trade Alert, which the headline excludes,
is added back from the October 2026 run (as of 5 October 2026, with the debiased strength
measure; 1{,}070 rows for 77 jurisdictions, EU members individually) with the weight the
composite gave it before (NIPO 4, IEA PAMS 3, CAT 2), IEA PAMS still required. NIPO is observed
in Upstream and Midstream cells only, in 841 of the Strategic Index cells (27 per cent). For
the jurisdictions it covers, it changes the policy term substantially: the within-country
Spearman correlation between the policy ranks with and without it has a median of 0.88 and
ranges from 0.00 (Taiwan) to 1.00. The Strategic Index absorbs much of this (pooled correlation
0.97, median within-country correlation 1.00), but the allied picks do not: 34 of the 57 picks
are unchanged, 27 of the 50 among members with a NIPO score, and the picks shift from Batteries
Midstream (27 to 15) towards Nuclear Upstream (8 to 14). A policy term built from announced
industrial-policy measures thus points to different sectors than one built from deployment
policies and climate ratings; the headline excludes it (Section~\ref{subsec:policy_nipo}).

\paragraph{Data choices of the final build and stability of the picks.} The last revision of
the inputs (11 and 12 October 2026) filled the fossil production that the Energy Institute does
not report from EIA International Energy Statistics, added Upstream trade baskets for Electric
Vehicles, Wind, Nuclear and Solar, rebuilt the energy-trade file from the all-reporter Comtrade
pull, and introduced the CM-FDI project data into Friendshore. Table~\ref{tab:robustness_data}
and Figure~\ref{fig:robustness_picks_data} report what these choices do. Setting the EIA-filled
production to zero, the earlier rule (the headline carries 807 EIA-sourced rows in the
import-dependence theme and 656 in the production theme), barely moves the pillars: the
within-pair correlations in Oil, Gas and Coal Upstream are between 0.97 and 0.98 in $\ES$ and
0.95 and 0.96 in $\EO$, and all 57 picks are unchanged. Removing the Trade category from the
four Upstream cells whose baskets are new moves those pairs more (median within-pair
correlations of 0.94 in $\ES$ and 0.73 in $\EO$, with a minimum of 0.58 in Solar Upstream $\EO$),
returns the Strategic Index to 192 countries, because the newly imputed Trade category had pushed
small economies over the rankability threshold, and changes 4 of the 57 picks, three of them to
Nuclear Upstream. Switching off the CM-FDI parts of Friendshore is not re-run on the final
outputs, because the dyad tables do not fit the memory budget of the replication; in the CM-FDI
integration study, on the inputs of 11 October 2026, the core-view PSI without them has a
Spearman correlation of 0.9998 with the CM-FDI build over 192 ranked partners, the largest rank
change is 5 places and the top 15 are the same, and switching off the Chinese-parent exposure
component or the CM-FDI investment ties alone gives correlations of 0.9999.

The revision of the energy-trade file could not be undone here, because the superseded file is
not kept, but the allied picks recorded at two interim builds show how much it mattered. On
8 October 2026, before the trade file was completed, 45 of the 56 members then in the coalition
had a pick, 17 of them Electric Grid Midstream and 11 Batteries Midstream; on 10 October, after
it, the corresponding counts were 4 and 26 of 57, and in the final build they are 4 and 27. The
shift is attributable mainly to the completed file, which filled the missing Electric Grid
Midstream codes. The category-drop
variants confirm that Trade is the input on which the picks rest most: without it 29 of the 54
members that keep a pick keep the baseline one, Electric Grid Midstream picks rise to 10 and
Batteries Midstream falls to 13, and measuring Trade on the Atlas basis instead keeps 49 of 57.
Technological Readiness, a technology-level constant that cannot reorder countries within a pair
but shifts a country's sectors against one another, comes next (32 of 57 unchanged; Electric
Grid Midstream would take 27 picks), followed by Technology Demand (41) and Foreign Dependency
(44). Taking both pillars from the 2025 vintage keeps 48 picks and from the 2020 vintage 44.
These comparisons say that the allied pick is an informative but fragile summary: it is stable
to the fill rules and the basket additions of the final revision, and sensitive to how trade is
measured and to which categories enter.

\begin{figure}[tbp]
  \centering
  \includegraphics[width=\textwidth]{robustness_picks_data}
  \caption{Stability of the allied picks under the data inputs. Each point is the number of
  coalition members, of 57, whose pick is the baseline one when the Strategic Index takes pillars
  rebuilt under a variant. Upper panel: one category dropped from the pillars that score it, the
  remaining categories renormalised (without Trade three members have no eligible rankable
  sector, so 29 of the 54 with a pick keep it). Lower panel: alternative inputs, namely both
  pillars from the 2020 or 2025 vintage build, the Trade category on the Atlas basis of the 2025
  vintage, the Trade category removed from the four Upstream cells whose trade baskets were added
  on 11 October 2026, and the fossil production that the Energy Institute does not report set to
  zero instead of the EIA value. Blue marks inputs revised in the final data fixes of
  11--12 October 2026. The Policy index and the endowment screen are the headline ones. The
  vertical line marks all 57 members. Source: category scores in \texttt{index\_outputs.rds},
  theme tables, \texttt{data/processed/vintages}; \texttt{scripts/paper/fig\_robustness.R}.}
  \label{fig:robustness_picks_data}
\end{figure}

% Generated by scripts/paper/fig_robustness.R; do not edit by hand.
\begin{table}[tbp]
  \centering
  \begin{threeparttable}
  \caption{Data and policy choices of the final build: rankable cells, rankings and allied picks}
  \label{tab:robustness_data}
  \footnotesize
  \setlength{\tabcolsep}{4pt}
  \begin{tabular}{lccccccc}
    \toprule
    & \multicolumn{2}{c}{Rankable cells} & \multicolumn{2}{c}{Median $\rho$, touched pairs} & \multicolumn{3}{c}{Strategic Index} \\
    \cmidrule(lr){2-3}\cmidrule(lr){4-5}\cmidrule(lr){6-8}
    Variant & $\ES$ & $\EO$ & $\ES$ & $\EO$ & Cells & Countries & Same pick \\
    \midrule
    Baseline & 4{,}163 & 4{,}301 & -- & -- & 3{,}154 & 174 & 57/57 \\
    \multicolumn{8}{l}{\emph{Data}}\\
    \quad Zero in place of the EIA fill & 4{,}164 & 4{,}301 & 0.98 & 0.96 & 3{,}155 & 174 & 57/57 \\
    \quad New Upstream baskets removed & 4{,}161 & 4{,}307 & 0.94 & 0.73 & 3{,}144 & 192 & 53/57 \\
    \multicolumn{8}{l}{\emph{Policy index}}\\
    \quad Both components required & = & = & -- & -- & 1{,}666 & 65 & 46/46 \\
    \quad IEA PAMS alone & = & = & -- & -- & 3{,}154 & 174 & 45/57 \\
    \quad October 2026 NIPO added back & = & = & -- & -- & 3{,}154 & 174 & 34/57 \\
    \multicolumn{8}{l}{\emph{Allied pick rule}}\\
    \quad Endowment screen off & = & = & -- & -- & 3{,}154 & 174 & 16/57 \\
    \bottomrule
  \end{tabular}
  \begin{tablenotes}[flushleft]\footnotesize
    \item \emph{Notes:} Rankable cells: headline cells of sovereign economies with at most half of the index weight imputed, after the variant (= unchanged). Median $\rho$: median across the touched technology--stage pairs (Oil, Gas and Coal Upstream for the EIA fill; Electric Vehicles, Wind, Nuclear and Solar Upstream for the baskets) of the within-pair Spearman correlation between the baseline and the variant over the baseline's rankable cells. Strategic Index: scored cells and countries, and the coalition members whose allied pick is the baseline one over the members with a pick. Zero in place of the EIA fill: the oil, gas and coal production that the Energy Institute does not report (807 EIA-sourced rows of the import-dependence theme and 656 of the production theme in the headline) set to zero, the rule until 11 October 2026, both themes rebuilt with the pipeline's builders. New Upstream baskets removed: the Trade category of the four Upstream cells whose trade baskets were added on 11 October 2026 dropped and both pillars rebuilt; the uranium codes moved out of Nuclear Midstream and the rare-earth codes moved out of EV and Wind Midstream are not restored. Policy rows as in Table~\ref{tab:robustness_picks}. The partnership side is not rebuilt: in the CM-FDI integration study (inputs of 11 October 2026), switching off both CM-FDI parts of Friendshore gives a core-view PSI with a Spearman correlation of 0.9998 with the CM-FDI build over 192 ranked partners, a largest rank change of 5 places and 15 of the top 15 in common. Source: \texttt{scripts/paper/fig\_robustness.R} (\texttt{data\_choices} and \texttt{policy} in \texttt{paper/results/robustness.json}).
  \end{tablenotes}
  \end{threeparttable}
\end{table}


\paragraph{Vertical complementarities.} The pillars score each stage of a technology's supply
chain separately. Supply shocks can propagate downstream along input--output links
\citep{acemoglu2012network, baqaee2019macroeconomic, carvalho2021supply}, and productive
capabilities in one activity can make related activities easier to enter
\citep{hidalgo2007product}, so a stage's opportunity might depend on the economy's capability in
the adjacent stage, and a stage's security on the security of its inputs. A companion study
tested both channels before any coupling of the stages was considered. For opportunity it used
the Atlas of Economic Complexity panel for 172 economies (1995--2023) and regressed the
five-year change in a stage's S-curved world market share, and entry into revealed comparative
advantage \citep{balassa1965trade}, on the adjacent stage's market-share index, with
country--year and technology--year fixed effects and standard errors clustered by country. The
cross-sectional association is positive and precisely estimated (0.046, standard error 0.007,
for Midstream on Upstream; 0.102, standard error 0.015, for the reverse), but it is a persistent
co-location rather than evidence of propagation: it is concentrated in fossil fuels and nuclear
power (0.062 and 0.151) and close to zero for the clean-technology chains (0.012 and 0.028, both
insignificant); with country--technology fixed effects it is a precisely estimated zero
($-0.001$, standard error 0.005, and 0.004, standard error 0.015); lags of the stage's own
capability absorb about a third of it (0.030 and 0.067); and it does not improve out-of-sample
prediction (the within-cell AUC for Midstream entry is 0.800 without the term and 0.794 with it,
and 0.739 and 0.741 for Upstream). The 2023 co-location lift, the ratio of the joint to the
independent probability of comparative advantage in both stages, is 4.4 for nuclear and 2.7 for
coal but 0.97 for wind and below one for gas, electric vehicles and the grid. For security the
study built a sourcing-weighted input-security measure, the domestic share times the economy's
own supply security plus the import-share-weighted security of its suppliers (half the
supplier's $\ES$ in the input stage, half the dyadic alignment). Its 2020 value does not predict
import shortfalls in 2021--22 (coefficient 0.10, standard error 0.31) or Midstream export growth
after 2020 (0.02, standard error 0.19), whereas the concentration of sourcing does (negative,
$p=0.035$); and the post-2022 shift of Upstream input trade towards aligned suppliers, 0.58 log
points (standard error 0.11) over all inputs, falls to 0.07 (standard error 0.06) once Russia
is excluded.

Table~\ref{tab:robustness_vertical} reports what the two terms would do to the indices if they
were added at the weights the evidence supports, each next to a placebo. The $\EO$ linkage term
is the adjacent stage's S-curved Atlas market share minus the economy's mean of that measure over
the nine basketed headline technologies, S-curved within the cell, so that it rewards an economy
for being relatively strong in the adjacent stage of this technology rather than for being a
large exporter overall; at the evidence weights (0.07 $\EO$ weight points in Midstream and 0.23
in Upstream, 0.6 and 1.5 per cent of a cell's mean observed weight) it leaves the within-pair
rankings practically unchanged (median correlation 0.999, minimum 0.987) and moves one pick,
Luxembourg's, whose margin is the smallest in the coalition; a placebo with the scores permuted
within each cell moves between none and one (5th to 95th percentile over 200 draws). The $\ES$ input-security term at a 15 per cent share moves the
rankings more (median within-pair correlation 0.92 in Midstream and 0.91 in Downstream, minimum
0.56 in Batteries Downstream) and changes 10 picks, but a placebo with the same scores permuted
within each cell changes 11 (5th to 95th percentile over 200 draws, 8 to 16), and a constant
placebo, which only rescales the other categories, changes 3. The term thus moves picks no more
than noise of the same size would. We therefore keep both pillars as uncoupled stage scores. The
evidence identifies co-location, not propagation; the weights it would justify are at most 1.5 per
cent of a cell's weight; and an adjacent-stage term in levels largely measures the overall size
of an economy's exports. Four extensions could change this conclusion and are left for future
work: plant-level co-location from the CM-FDI and BNEF project data; input shares weighted by
input--output coefficients rather than trade baskets; a test of the sourcing-concentration
measure on the 2023 export controls on graphite and gallium; and a co-location term in the
network design of Section~\ref{sec:network} rather than in the node scores.

\paragraph{Partnership index: the viewpoint.} The headline PSI averages the dyadic index over the
viewpoints of the 31 core members (equation~\eqref{eq:psi}); the full
coalition (57 reporters) and the world (218 reporters) are alternatives
(Figure~\ref{fig:robustness_psi_views}). With the rankability rule applied to the PSI, the
core and full-coalition views rank 192 of 219 partners and the world view 188. The full
coalition gives practically the same ranking as the core (Spearman correlation 0.98 over the 192
partners ranked in both, with 13 of the top 15 in common); its top 15 adds China and India and
drops Ireland and Sweden, and China moves from 44th to 9th, a move consistent with the higher
dyadic alignment between China and several of the partners the full coalition adds. The world view
differs more (0.88 over 188 partners, with 10 of the top 15 in common, and a median absolute rank
change of 17 places): China, India, Indonesia, Viet Nam and Malaysia enter its top 15, and
Belgium, the Czech Republic, Ireland, Poland and Sweden leave it. The full-coalition and world
views correlate at 0.95. The choice of viewpoint is thus a choice about whose interests the index
represents rather than a technical detail; the core view is the headline because the network
design and the allied picks are built for the coalition, and Section~\ref{sec:bilateral}
reports both alternatives. The rankability threshold matters more for the world view than for
the core: at $\tau=0.3$ the world view ranks 102 partners and the core view 159.

\begin{figure}[tbp]
  \centering
  \includegraphics[width=\textwidth]{robustness_psi_views}
  \caption{The partnership index from three viewpoints. Each point is a partner ranked in both
  views (rankable under $\tau=0.5$), placed by its percentile position in the core-coalition
  view, the headline average over the 31 core members' ratings (horizontal), and in an
  alternative view (vertical): left, the full coalition of 57 reporters; right, the world of
  218 reporters. Blue marks the top 15 of the core view; the labelled points are the three of
  them with the lowest position in the alternative view. The panel titles give the Spearman rank
  correlation with the core view and the number of core-view top-15 partners that are also in the
  alternative top 15. The line is the 45-degree line. Source: \texttt{psi\_tbl.rds},
  \texttt{psi\_full\_coalition\_tbl.rds}, \texttt{psi\_world\_tbl.rds};
  \texttt{scripts/paper/fig\_robustness.R}.}
  \label{fig:robustness_psi_views}
\end{figure}

\paragraph{What the analysis does and does not establish.} Within the ranges explored, the
within-pair rankings of $\ES$ are robust to the category weights, the curvature and, over the
rankable cells, the imputation rules, and those of $\EO$ to the weights and the curvature but
not to the imputation rules; the Strategic Index orders a country's sectors stably, while the
allied pick, a single arg max, is robust only where its margin is large, and it turns on how
trade is measured, on the rankability threshold, on the composition of the policy term and on
the endowment screen. Several limits qualify this. First, there is no external validation of the
headline indices: we have no outcome series, such as realised supply disruptions for $\ES$ or
subsequent export or investment growth for $\EO$, against which to test whether the pillars
measure what they are meant to measure; the only outcome tests in the paper are those of the
vertical-complementarity study, which concern candidate terms rather than the pillars. Stability
shows that the rankings are determined by the data and the construction rather than by fine
choices of parameters, not that they are right \citep{ravallion2012mashup}. Second, much of the
stability is mechanical. Within a pair only a few categories vary, neither a constant category
nor a monotone re-curving of a single varying component can reorder countries, and the cell
constants (Energy Prices, the world-concentration terms, technological readiness) are inert by
construction; robustness of this kind says little about the constant categories, which matter
for comparisons across cells and hence for the Strategic Index. Third, the perturbations are
local: the Dirichlet draws are centred on the configured weights and vary them by about 30 per
cent, and the other experiments vary one element at a time. Fourth, the vintage comparison
covers only the time-varying sources, which carry about a third of the average cell's weight, so
it cannot say how stable the full index is over time. Fifth, some data choices are not varied
here: the exclusion of electricity from the grid basket, the exclusion of Hong Kong's
re-exports from the Comtrade export measures, the treatment of non-reporting economies through
their partners' mirror declarations, the rank compression of the EU demand weight, the regional
IEA shares that every EU member inherits in the foreign-dependency and market-share variables,
and the zero counted for unquantified USGS holdings; and the CM-FDI comparison rests on the
inputs of 11 October 2026 rather than on the final run. Finally, the sensitivity of the network
design to its own settings, including the calibration of the diversification weight $\lambda$,
which sits on a knife edge, is reported with the design in Section~\ref{sec:network}.

% Appendix sensitivity table of the vertical-complementarity study (user decision 2026-10-12:
% paper only, no index change). The integrating section may move this \input to an appendix.
% Generated by scripts/paper/fig_robustness.R; do not edit by hand.
\begin{table}[tbp]
  \centering
  \begin{threeparttable}
  \caption{Vertical complementarities: sensitivity of the pillars and the allied picks to an adjacent-stage term}
  \label{tab:robustness_vertical}
  \footnotesize
  \setlength{\tabcolsep}{4pt}
  \begin{tabular}{llccc}
    \toprule
    Term & Weight & \makecell{Median within-\\pair $\rho$} & \makecell{Minimum\\$\rho$} & \makecell{Picks\\changed} \\
    \midrule
    $\EO$ linkage, within-country deviation & Mid 0.07, Up 0.23 & 0.999 & 0.987 & 1 \\
    \quad Placebo: scores permuted within cell & same & 0.999 [0.999, 0.999] & -- & 0 [0, 1] \\
    \addlinespace
    $\ES$ input security, sourcing-weighted & 15\% share & 0.917 & 0.556 & 10 \\
    \quad Placebo: constant 0.5 & same & 1.000 & 1.000 & 3 \\
    \quad Placebo: scores permuted within cell & same & 0.902 [0.891, 0.909] & -- & 11 [8, 16] \\
    \bottomrule
  \end{tabular}
  \begin{tablenotes}[flushleft]\footnotesize
    \item \emph{Notes:} $\EO$ linkage: a category added to the Midstream and Upstream $\EO$ cells, the S-curved world market share of the adjacent stage (Atlas of Economic Complexity, HS92 baskets, 2023) minus the economy's mean of that measure over the 9 basketed headline technologies, S-curved within the cell (Midstream 1{,}890; Upstream 1{,}890 cells). Weights in $\EO$ weight points (the nominal category weights are 2 to 4): the evidence weights of the vertical-complementarity study, 0.6 times the ratio of the adjacent-stage to the feasibility coefficient in the five-year market-share regression, upper bounds because about a third of the coefficient is absorbed by lags of own capability. $\ES$ input security: the domestic supply share times the economy's own supply security plus the import-share-weighted security of its suppliers (half the supplier's $\ES$ in the input stage, half the dyadic alignment), 2024 sourcing, S-curved within the cell and given 15 per cent of the Midstream and Downstream cells' weight (Downstream 1{,}626; Midstream 1{,}893 cells). Median and minimum $\rho$: within-pair Spearman correlation with the baseline pillar over the baseline's rankable cells of the receiving stages; for the permutation placebo, the median and the 5th and 95th percentiles over 200 draws of the per-draw median. Picks changed: coalition members (of 57) whose allied pick differs from the baseline when the Strategic Index takes the modified pillar (for the placebo, median and 5th--95th percentiles over the draws). Neither term enters the headline indices. Source: \texttt{scripts/paper/fig\_robustness.R} (\texttt{vertical} in \texttt{paper/results/robustness.json}), from the study's inputs.
  \end{tablenotes}
  \end{threeparttable}
\end{table}

